% ch7_d §7.10–7.11 (书页 239–245 / p253–p259)
%--C-TAIL--
%JOIN%也应当有类似的结果，只是改由 $\zeta_{h}$ 扮演“每个载流子所携带的热量”的角色，并且此时电荷为正。当两类载流子同时存在时，我们得到的将是一个平均的温差电功率（thermo-electric power），介于电子的负贡献与空穴的正贡献之间而取得平衡，各自按其在总电流中所占的份额加权。

有趣的是，不妨想一想：对于金属，在考虑“空穴球”时，我们有哪两种可供选择的处理办法。如同 §6.6，我们可以把它看作一个电子的 费米面，其 $e$ 为负，但具有一个奇特的性质：电子的面积、速度和弛豫时间全都随能量的增加而\emph{减小}，于是在式 (7.104) 中 $\partial\ln\sigma(\mathcal{E})/\partial\mathcal{E}$ 也是负的，从而 $Q$ 为正。更自然的做法是把它看作一个由“空穴”组成的球，空穴的数目和迁移率随其“能量”自中心向外而\emph{增加}，但它们携带的是正电荷。两种办法的结果相同。

%FIG{130}: {计算温差电功率的两种可供选择的模型：（a）电子 费米面；（b）“空穴”费米面。}

\section{晶格热传导}

在电绝缘体和半导体中，仍可能有经由格波的良好的热传导。这是一个极为复杂的课题，我们只对它作一点粗浅的研究。

设想我们把声子当作一种粒子气体来处理。初等分子运动论中已经证明，热导率为
\begin{equation*}
\kappa=\tfrac{1}{3}\,Cv\Lambda,\tag{7.110}
\end{equation*}
其中 $C$ 是比热，$v$ 是速度，$\Lambda$ 是载流子的平均自由程（mean free path）。

在高温下，我们可以认为声子具有其通常的比热 $3Nk$，如式 (2.60)；也可以认为它们以声速 $s$ 行进。但要计算 $\Lambda$，我们还须作某些粗略的近似。其中最简单的一种，与 §7.5 中估计电子平均自由程时所用的论证相仿。

让我们回忆一下，晶格的膨胀 $\Delta$ 会局部地改变声速，相对改变量为
\begin{equation*}
\frac{\delta s}{s}=\gamma\Delta,\tag{7.111}
\end{equation*}
其中 $\gamma$ 是 Grüneisen 常数。这个常数在式 (2.120) 中讨论热膨胀理论时已经定义过。

密度的局域热涨落将引起散射，其强度正比于这个量的均方。与式 (7.59) 直接比较，可以写出
\begin{equation*}
\frac{1}{\Lambda}\sim\frac{\gamma^{2}\overline{\Delta^{2}}}{a},\tag{7.112}
\end{equation*}
亦即
\begin{equation*}
\Lambda\sim\frac{Ds^{2}}{NkT}\,\frac{a}{\gamma^{2}},\tag{7.113}
\end{equation*}
其中 $a$ 是具有晶格常数量级的一个长度。于是，利用式 (7.110) 并略去数值因子，便得到
\begin{equation*}
\kappa\sim\frac{Ds^{3}}{\gamma^{2}T}\,a,\tag{7.114}
\end{equation*}
其中 $D$ 是密度。这个公式有意思之处在于，除了 $a$ 以外，式中出现的全是宏观量。

改写式 (7.113) 的另一种方式是用熔化温度来表示：如式 (7.62)，我们有
\begin{equation*}
\Lambda\sim\frac{20}{\gamma^{2}}\,\frac{T_{m}}{T}\,a.\tag{7.115}
\end{equation*}

这些公式在函数行为和数量级上是合理的；晶格热导率对绝对温度倒数的依赖关系已被很好地理解，尽管其余因子是相当粗糙的。

然而，这只在高于 Debye 温度的高温下才成立，而且实际情形远不像表面上那样简单。事实是，“密度的热涨落”与“携带热流的声子”并无本质区别，整个系统应当像 §2.10 中那样用声子--声子相互作用来分析。

于是我们得到一个非常特别的结果。设想只有晶体动量守恒的声子--声子 N 过程。设想我们建立起一个在模式 $\mathbf{q}$ 中含有 $n_{\mathbf{q}}$ 个声子的态——即一个具有净晶体动量
\begin{equation*}
\mathbf{P}=\sum_{\mathbf{q}}n_{\mathbf{q}}\hbar\mathbf{q}\tag{7.116}
\end{equation*}
的态。这时，声子--声子 N 过程并不改变 $\mathbf{P}$ 的值。因此，与 $\mathbf{P}$ 相联系的任何净热流都不会耗散：热导率必须为无穷大。

实际上，我们得到的是声子沿样品的\emph{对流}（convection）。正如在\emph{流动的气体}中那样，粒子之间的碰撞允许动量发生交换，并在气体内部建立起某种程度的局域平衡；除了通过与器壁的相互作用而间接起作用之外，它们并不会减慢热的质量输运。当然，对于气体中的\emph{热传导}，条件就完全不同了；我们刻意坚持不许粒子有任何质量输运，这时碰撞才变得重要起来，因为朝一个方向运动的“热”粒子把能量交给朝相反方向运动的“冷”粒子。但在声子气体中，“粒子”并没有净的守恒性，所以我们不能强加这样的条件。携带热量所需的额外声子在热端产生，流向冷端，并在那里被消灭。

晶体热导率之所以有限，是由于 U 过程（U-process）的存在。由于在这类过程中晶体动量不守恒，量 (7.116) 不能保持恒定，热流于是被耗散。这种情形下热阻的计算相当复杂，但其结果在高温下正与式 (7.114) 和 (7.115) 相合。这个非常朴素的近似方法结果证明是有道理的。

但在低温下，声子--声子 U 过程引起的电阻迅速降为零。我们回忆一下（§2.11）波矢所满足的几何条件
\begin{equation*}
\mathbf{q}+\mathbf{q}'=\mathbf{q}''+\mathbf{g},\tag{7.117}
\end{equation*}
以及频率所满足的能量条件
\begin{equation*}
\nu+\nu'=\nu''.\tag{7.118}
\end{equation*}
这样，$\mathbf{q}''$ 必须足够大，使其能量等于 $\mathbf{q}$ 与 $\mathbf{q}'$ 的能量之和；但要发生 U 过程，这些矢量之和必须越出第一 Brillouin 区，因此 $\mathbf{q}''$ 的长度不可能比 $\frac{1}{2}\mathbf{g}$ 小很多。在 Debye 模型中，这意味着
%FIG{131}: {声子--声子 U 过程。}
\begin{equation*}
q''\geqslant\beta q_{D},\tag{7.119}
\end{equation*}
其中 $\beta$ 是一个分数，量级为 $\frac{1}{2}$ 或 $\frac{2}{3}$。换句话说，
\begin{equation*}
\hbar\nu''\geqslant\beta k\Theta.\tag{7.120}
\end{equation*}
形如式 (7.118) 的 U 过程出现的概率，正比于相应声子模式占据数的乘积
\begin{align*}
n_{\mathbf{q}}\,n_{\mathbf{q}'}&\approx\exp(-\hbar\nu/kT)\exp(-\hbar\nu'/kT)\\
&\approx\exp(-\hbar\nu''/kT)\\
&\approx\exp(-\beta\Theta/kT)\tag{7.121}
\end{align*}
在低温下它迅速趋于零。我们说倒逆散射被“冻结”了，热导率按
\begin{equation*}
\kappa\sim T^{n}\exp\left(\frac{\beta\Theta}{T}\right),\tag{7.122}
\end{equation*}
趋于无穷，其中 $n$ 是一个依赖于模型细节的指数。

实际发生的情形是：声子的平均自由程变得可以与晶体的尺寸相比拟。设想我们处理的是一根直径为 $D$ 的棒。于是在式 (7.110) 中可以取
\begin{equation*}
\Lambda=D,\tag{7.123}
\end{equation*}
从而
\begin{equation*}
\kappa\propto T^{3}D,\tag{7.124}
\end{equation*}
这是因为低温下比热服从 $T^{3}$ 定律（§2.4）。这个 $T^{3}$ \emph{尺寸效应}（size effect）在实验上有确凿的证据。

然而，还存在使声子受到散射的其他机制。即使在“完美”晶体中，原子也并非全都等价；对每一种化学元素来说，通常都是一些不同质量的同位素的混合。这些质量上的差别能够散射声子。我们很容易——比如用初等弹性理论——证明 \emph{Rayleigh 公式}：密度为 $D$ 的介质中一个点质量 $\delta M$ 对波数为 $q$ 的波的散射截面为
\begin{equation*}
\sigma(q)=\frac{q^{4}}{4\pi D^{2}}\,(\delta M)^{2}.\tag{7.125}
\end{equation*}
倘若晶体中每个原子的质量都可能偏离平均值这么多，那么声子的平均自由程将为
\begin{align*}
\Lambda(q)&=\frac{4\pi NM^{2}}{q^{4}(\delta M)^{2}}\\
&=\left(\frac{M}{\delta M}\right)^{2}\frac{4\pi a}{(aq)^{4}},\tag{7.126}
\end{align*}
其中 $N=1/a^{3}$。

$\Lambda(q)$ 随 $q$ 的强烈变化现在成了一个难题。自然的做法是把式 (7.110) 加以推广，写成
\begin{equation*}
\kappa=\frac{1}{3}\int C(q)\,v(q)\,\Lambda(q)\,\mathrm{d}\mathbf{q}\tag{7.127}
\end{equation*}
把各种模式的贡献加起来。但这会在 $q\to 0$ 处导致一个奇异性：在那里比热 $C(q)$ 和速度 $v(q)$ 实际上是常数，而 $\Lambda(q)\propto q^{-4}$；于是，对长波模式的分布利用类似于 (2.56) 的 Debye 谱，便得到
\begin{equation*}
\kappa\propto\int\frac{q^{2}\,\mathrm{d}q}{q^{4}}\sim\int_{0}\frac{\mathrm{d}q}{q^{2}}\to\infty\tag{7.128}
\end{equation*}

这一困难的解决是一个复杂的问题，但论证的实质是：即使在低温下，晶体中也总有强烈的声子--声子 N 过程在进行。这些过程本身并不产生热阻，但它们倾向于使各种模式彼此趋于平衡，因而不允许那些不受同位素强烈散射的长波模式携带过大的热流。

倘若声子--声子 N 过程确实非常强烈，我们就可以论证：式 (7.127) 中应当加以平均的量是 $1/\Lambda(q)$，这样积分才是有限的。在这种情形下，我们可以这样来估计用于分子运动论公式 (7.110) 中的平均自由程：假定在温度 $T$ 下只有一个“典型”的声子模式，其波矢为
\begin{equation*}
\overline{q}\sim\frac{T}{\Theta}\,q_{D}.\tag{7.129}
\end{equation*}
于是我们应当写出
\begin{equation*}
\Lambda\sim\Lambda(\overline{q})\propto T^{-4},\tag{7.130}
\end{equation*}
这将给出
\begin{equation*}
\kappa\propto\left(\frac{M}{\delta M}\right)^{2}\frac{1}{T}.\tag{7.131}
\end{equation*}
然而，更细致的分析——采用声子--声子相互作用的唯象模型——给出一种更复杂的行为：
\begin{equation*}
\kappa\propto\frac{M}{\delta M}\,\frac{1}{T^{1/2}}.\tag{7.132}
\end{equation*}
这与实验观测大致相符。

\section{声子曳引}

金属中的电子受声子散射。反过来，声子也受电子散射。一般说来，这种散射十分强烈，以致与电子气的热导率相比，晶格热传导就很难观察到了。

然而，在金属和半导体中，有一个由声子引起的重要电学效应。设想一根导线中流有电流。如式 (7.27) 所示，这可以描述为 费米面在 $\mathbf{k}$ 空间中的位移，位移量为 $(e\tau/\hbar)\mathbf{E}$。这个电子系统正在与声子相互作用。如果只有电子--声子相互作用，声子系统将力图与\emph{位移了的}电子系统达到平衡；也就是说，声子系统本身将倾向于在自己的动量空间中发生位移。

但声子系统的这样一种位移，对应于声子在真实空间中的“漂移”——也就是说，它对应于一个热流。设想声子的漂移速度与电子的漂移速度 (7.31) 相同。若 $C_{L}$ 是整个声子系统的比热，就会存在一个晶格热流
\begin{equation*}
\mathbf{U}_{L}\sim C_{L}T\,\delta\mathbf{v}.\tag{7.133}
\end{equation*}
但它又如式 (7.32) 那样伴随有一个电流。于是将有一个 Peltier 系数（参见式 (7.102)）
\begin{equation*}
\frac{U_{L}}{J}\sim\frac{C_{L}T}{ne},\tag{7.134}
\end{equation*}
亦即一个\emph{晶格温差电功率}（lattice thermo-electric power）
\begin{equation*}
Q_{L}\sim\frac{C_{L}}{ne}.\tag{7.135}
\end{equation*}

与普通的电子贡献 (7.106) 相比，这将是一个很大的效应；在后者中起作用的是电子的小得多的比热——如式 (7.107)——而不是晶格的比热。但当然，我们原本曾假定：除了电子--声子相互作用之外，声子不受任何其他机制的散射。在寻常温度下情况并非如此；那里存在着 §7.10 中讨论过的声子--声子 U 过程等等，它们使这个 \emph{Gurevitch 效应}（Gurevitch effect）大为减弱。

然而，在低温下，当声子--声子相互作用变得可以忽略时，\emph{声子曳引}（phonon drag）效应就很容易探测到了。在这种情况下，晶格比热按 $T^{3}$ 变化。对于单价金属，公式应当是
\begin{equation*}
Q_{L}\sim\frac{k}{e}\,\frac{4\pi^{4}}{5}\left(\frac{T}{\Theta}\right)^{3},\tag{7.136}
\end{equation*}
这是一个很大的\emph{负}温差电功率，其符号由电子电荷 $e$ 的符号决定。

但这并不总是实际观察到的情形，理由如下。“声子分布以电子的漂移速度一道被曳引”这一假设，依赖于电子--声子过程中晶体动量的守恒。考虑（图 132(a)）一个普通的 N 过程：电子由态 $\mathbf{k}$ 被散射到态 $\mathbf{k}'$，同时发射出一个声子。
%FIG{132}: {电子对声子的曳引：（a）N 过程；（b）U 过程。}
如果这一过程恰好使 $\mathbf{k}$ 沿场方向的分量反转，那么声子便沿该方向发射。于是，使电流减慢的这个过程，就产生了沿载流子运动方向运动的过剩声子。

现在考虑在电子--声子 U 过程中会发生什么。如图 132(b) 所示，声子倾向于沿与载流子电流\emph{相反}的方向被发射；矢量 $\mathbf{k}$、$\mathbf{k}'$ 与 $\mathbf{q}$ 之和必须给出大的倒格矢 $\mathbf{g}$。这样，声子分布就不随电子一道被“曳引”；它反而可能向相反的方向漂移。

因此，$Q_{L}$ 的实际符号可正可负，取决于在金属的“理想”电阻中 N 过程与 U 过程哪一种更为重要。这一效应显然非常敏感
% 注：本文件末句在原书 p259 末行处中断（sensitive to...），由 ch7_e（p260 起）直接续接，勿加标点。
